% ============================================================
% Project : THE PHILOSOPHY OF VIABILITY
% Book    : Book I — Ontology of Viability
% Book ID : VS.B01
% File    : section_12_formal_reconstruction_of_invariance.tex
% Section : Chapter 07, Section 12
% Purpose : To reconstruct invariance formally by specifying
%           domains, transformations, respects, invariant
%           structures, and the limits of formal inference
%           concerning ontological invariance.
% ============================================================

\section{Formal Reconstruction of Invariance}
\label{sec:formal-reconstruction-of-invariance}

\subsection{The Formal Problem of Invariance}

Invariance concerns the preservation of a specified
feature, relation, structure, or condition under a
specified transformation or class of transformations.

Its formal reconstruction requires more than the
assertion that something remains unchanged. It
requires the identification of what is preserved,
under which transformation, within which domain,
and according to which criterion of preservation.

An invariance claim is therefore relational in its
formal structure. It connects an object or structure
to a transformation and specifies a respect in
which the transformation leaves something invariant.

The central thesis of this section is:

\begin{quote}
An invariance claim is formally meaningful only when
its domain, transformation, respect of preservation,
and criterion of invariance are sufficiently
specified.
\end{quote}

This thesis provides the basis for distinguishing
different forms of invariance and for preventing
formal invariance from being treated as an automatic
proof of ontological identity, persistence, continuity,
or stability.

The formal reconstruction must also distinguish
between the invariance of a mathematical expression
and the invariance of an ontological feature. The
former is determined by a formal system and its
interpretation. The latter requires an independently
justified account of the relevant ontological
structure.

\subsection{The Components of an Invariance Claim}

A complete invariance claim should identify at least
five components:

\begin{enumerate}
    \item \textbf{Domain:} the objects, structures,
    states, or relations to which the claim applies;

    \item \textbf{Transformation:} the operation or
    class of operations under which preservation
    is examined;

    \item \textbf{Respect:} the particular feature,
    relation, or aspect considered invariant;

    \item \textbf{Criterion:} the condition by which
    preservation is determined;

    \item \textbf{Scope:} the limits within which the
    claim is asserted.
\end{enumerate}

These components are not interchangeable.

A transformation may preserve one feature while
altering another. A structure may remain invariant
under one class of transformations but fail to
remain invariant under a broader class. A property
may be preserved within a restricted domain without
being preserved across every possible domain.

Consequently, the unqualified expression
``the structure is invariant'' is insufficient
unless the relevant conditions are already
established by context.

\subsection{A General Formal Schema}

Let:

\[
    D
\]

denote a specified domain, and let:

\[
    T : D \rightarrow D
\]

be a transformation defined on that domain.

Let \(R\) denote a specified respect of preservation,
and let:

\[
    \Phi_R : D \rightarrow V_R
\]

represent a formal characterization of the feature
being examined, where \(V_R\) is an appropriate
codomain.

A basic invariance condition may then be expressed as:

\[
    \forall x \in D,\qquad
    \Phi_R(T(x)) = \Phi_R(x).
\]

This expression states that the value assigned by
\(\Phi_R\) is preserved under \(T\).

The notation is schematic. It assumes that:

\begin{itemize}
    \item \(T\) is defined for the relevant elements
    of \(D\);

    \item \(\Phi_R\) is defined on the relevant
    transformed and untransformed elements;

    \item equality in \(V_R\) is an appropriate
    criterion for identifying preservation;

    \item the chosen representation captures the
    intended respect \(R\).
\end{itemize}

If the transformation is partial, the formula must
be restricted to its domain of definition. If the
transformation maps between different domains,
the schema must be adjusted accordingly.

More generally, let:

\[
    T : D_1 \rightarrow D_2.
\]

A preservation claim may then take the form:

\[
    \forall x \in D_1,\qquad
    \Phi_R^{(2)}(T(x))
    =
    \Phi_R^{(1)}(x),
\]

provided that the two mappings express comparable
features and that the relevant equality relation
has been specified.

This generalized form makes explicit that invariance
does not always require a transformation from a
domain into itself.

\subsection{Invariance Under a Class of Transformations}

An invariance claim may concern a single
transformation or a collection of transformations.

Let:

\[
    \mathcal{T}
\]

be a specified class of admissible transformations.

A feature is invariant under this class when:

\[
    \forall T \in \mathcal{T}\;
    \forall x \in D_T,\qquad
    \Phi_R(T(x)) = \Phi_R(x),
\]

where \(D_T\) denotes the relevant domain of
definition for each transformation.

The strength of the claim depends on the
transformation class.

If:

\[
    \mathcal{T}_1 \subseteq \mathcal{T}_2,
\]

then invariance under \(\mathcal{T}_2\) entails
invariance under \(\mathcal{T}_1\), provided that
the same domain, respect, and preservation criterion
are maintained.

The converse does not generally hold.

A feature may remain invariant under a restricted
class of transformations while changing under
a broader class.

This has an important methodological consequence:
the strength of an invariance claim cannot be
evaluated independently of the scope of its
transformation class.

\subsection{Transformation and Preservation}

The notion of transformation must be specified
according to the formal subject under examination.

A transformation may represent:

\begin{enumerate}
    \item a change in the representation of an object;
    \item a change in its formal coordinates;
    \item a change in its internal configuration;
    \item a change in its relational arrangement;
    \item a transition between formally described states;
    \item an operation on a mathematical structure;
    \item a change in the conditions under which
    an ontological feature is examined.
\end{enumerate}

These cases do not necessarily have the same
ontological significance.

For example, a transformation that changes only
the representation of a structure may preserve
the structure itself. A transformation that changes
the internal constitution of an existent may
preserve some relational or organizational features
while altering others.

The formal system must therefore distinguish
transformations that are merely representational
from transformations that are intended to describe
changes in the subject matter.

The preservation of a feature under a
representational transformation cannot automatically
be interpreted as its preservation under an
ontological transformation.

\subsection{Different Respects of Invariance}

The respect of invariance identifies what is
preserved.

A formal reconstruction may distinguish several
possible respects, depending on the subject matter:

\begin{itemize}
    \item \textbf{Numerical invariance:} preservation
    of a specified numerical value;

    \item \textbf{Structural invariance:} preservation
    of a specified structural arrangement;

    \item \textbf{Relational invariance:} preservation
    of a specified relation or relational pattern;

    \item \textbf{Functional invariance:} preservation
    of a specified functional relation or mapping;

    \item \textbf{Geometrical invariance:} preservation
    of a geometrical feature under a specified
    transformation;

    \item \textbf{Formal invariance:} preservation
    of a specified property within a formal system;

    \item \textbf{Ontological invariance:} preservation
    of a specified ontological feature under a
    transformation whose ontological relevance
    has been justified.
\end{itemize}

This classification is not intended as an exhaustive
taxonomy. It identifies different possible respects
whose meanings and relations must be established
in each application.

In particular, the use of the same word
``invariance'' across different domains does not
prove that all these forms share one ontological
foundation.

\subsection{Invariant Functions and Invariant Relations}

The preservation of a function and the preservation
of a relation require different formal treatment.

For a function:

\[
    f : D \rightarrow V,
\]

invariance under \(T\) may be expressed as:

\[
    f(T(x)) = f(x).
\]

For a relation:

\[
    R \subseteq D^n,
\]

preservation may be expressed by requiring that
the relation hold for a tuple if and only if it
holds for the transformed tuple, when the
transformation acts appropriately on each component:

\[
    R(x_1,\ldots,x_n)
    \leftrightarrow
    R(T(x_1),\ldots,T(x_n)).
\]

The biconditional expresses preservation in both
directions under the specified transformation.

A one-directional condition may instead be
appropriate when only preservation from the
original structure to the transformed structure
is required.

The distinction matters because a transformation
may preserve a relation without reflecting it.
Preservation and reflection are separate formal
conditions and should not be conflated.

Similarly, the invariance of a function's value
does not establish the invariance of every relation
involving that function.

The exact object of preservation must remain
explicit.

\subsection{Invariance and Group Actions}

In some formal settings, transformations form a
group or act on a domain through a group action.

Let \(G\) be a group acting on \(D\). The action
may be represented as:

\[
    \alpha : G \times D \rightarrow D.
\]

A function \(f\) is invariant under the action
when:

\[
    f(\alpha(g,x)) = f(x)
\]

for every admissible \(g \in G\) and \(x \in D\).

This formal structure is useful for studying
symmetries and other transformation-based
preservation conditions.

However, group structure is not a universal
requirement for invariance.

Some transformations do not form groups. They
may be irreversible, partial, non-invertible,
or organized by weaker algebraic structures.

Consequently, a formal reconstruction of
invariance must not assume that every relevant
transformation is reversible or that every
invariance relation is a symmetry.

The existence of a group action is a specific
formal condition, not a general ontological
presupposition.

\subsection{Invariance, Identity, and the Principles
of Being}

The relation between invariance and identity
requires particular care.

An invariance claim concerns the preservation
of a specified feature under a transformation.
An identity claim concerns the identity of an
existent.

The two may be related in a particular argument,
but neither is reducible to the other.

Suppose a feature \(\Phi_R\) is invariant under
a transformation \(T\):

\[
    \Phi_R(T(x)) = \Phi_R(x).
\]

This establishes the preservation of the specified
feature under the specified transformation.

It does not establish, without further premises,
that:

\[
    T(x) = x.
\]

Nor does it establish that the identity of the
existent is exhausted by \(\Phi_R\).

A transformation may alter many features while
preserving one. Conversely, a feature may change
without this fact alone settling whether the
existent retains its identity.

The formal reconstruction must therefore avoid
the following invalid transitions:

\[
    \text{Invariance}
    \;\not\Rightarrow\;
    \text{Identity},
\]

\[
    \text{Non-invariance}
    \;\not\Rightarrow\;
    \text{Loss of identity}.
\]

These expressions indicate the absence of an
automatic entailment, not the impossibility of
establishing a relation through additional
arguments.

The distinction also protects the three
foundational principles concerning being and
identity.

Analytical Distinction does not imply that
invariance is the foundation of identity.
Ontological Inseparability does not imply that
identity is defined by the preservation of
invariant features. Absence of Ontological
Priority does not establish that invariance
has no explanatory or formal relevance to
questions concerning identity.

Any proposed connection requires its own
argument and must specify the respect in which
the relation is asserted.

\subsection{Invariance and Persistence}

Persistence concerns the continued existence
of an existent across some specified interval
or sequence of conditions.

Invariance concerns the preservation of a
specified feature under a transformation.

The two may intersect, but their formal structures
are different.

A persistence claim may be represented
schematically as:

\[
    P(x,t_1,t_2),
\]

where \(P\) denotes a specified persistence
relation over an interval.

An invariance claim may be represented as:

\[
    \Phi_R(T(x)) = \Phi_R(x).
\]

Neither expression entails the other without
additional premises connecting the relevant
temporal and transformational structures.

An existent may persist while undergoing changes
that alter a particular feature. A feature may
remain invariant under a formal transformation
without the model establishing persistence
through time.

Consequently, persistence and invariance must
be reconstructed independently before any
connection between them is asserted.

\subsection{Invariance and Continuity}

Continuity may concern the organization of a
transition, the relation between successive
states, or the absence of specified forms of
disruption.

Invariance concerns preservation under a
transformation.

A continuous process need not preserve every
feature. A discontinuous transition may preserve
a particular formal or relational feature.

The two notions therefore require distinct
formal criteria.

For example, continuity may be represented
through a topological or analytical structure,
whereas invariance may be represented through
a transformation and a preservation condition.

The choice of mathematical representation
depends on the subject matter. No single formal
schema should be assumed to capture every
ontological use of continuity.

The formal reconstruction must avoid treating
continuity as a consequence of invariance or
invariance as a consequence of continuity
unless the relevant assumptions establish
that relation.

\subsection{Invariance and Stability}

Stability concerns the behavior of a system
under specified disturbances, perturbations,
or changes in its conditions.

Invariance concerns the preservation of a
specified feature under specified transformations.

A system may be stable in a defined sense while
some of its features vary. It may also preserve
a feature under a transformation without
satisfying a broader stability condition.

A stability claim typically requires additional
formal resources, such as:

\begin{itemize}
    \item a state space;
    \item a class of perturbations;
    \item a criterion of acceptable deviation;
    \item a measure of response or convergence;
    \item a specified temporal or iterative structure.
\end{itemize}

These requirements are not automatically present
in an invariance claim.

The reconstruction must therefore avoid deriving
stability from invariance alone or using the two
terms as interchangeable descriptions of
preservation.

\subsection{Formal Invariance and Ontological Invariance}

Formal invariance is established within a
specified formal system and interpretation.

Ontological invariance concerns the preservation
of an ontological feature under a transformation
whose relevance to the intended subject matter
has been justified.

The distinction may be represented schematically:

\[
    \text{Formal Invariance}
    \neq
    \text{Ontological Invariance}.
\]

The inequality indicates that the two claims
have different justificatory requirements.

A formal invariant may be an artifact of the
chosen representation. For example, a quantity
may remain constant because the formalism
defines the transformation in a way that
preserves it.

To establish ontological significance, the
reconstruction must explain why the represented
quantity corresponds to a genuine feature of
the intended subject matter and why the
transformation captures a relevant ontological
change.

This requires a justified interpretive bridge
between the formal model and the ontology.

The bridge cannot be supplied by notation
alone.

\subsection{Invariance Across Alternative Models}

The model-based method developed in Section 11
provides a further test for invariance claims.

Let \(\mathcal{M}\) be a specified class of
admissible models.

An invariance claim may be examined across
this class by asking whether the relevant
preservation condition holds in every model:

\[
    \forall \mathfrak{M} \in \mathcal{M},
    \qquad
    \mathfrak{M} \models \operatorname{Inv}(R,T).
\]

Here, \(\operatorname{Inv}(R,T)\) denotes
the formal expression of invariance in the
specified respect under transformation \(T\).

If the condition holds in every admissible
model, the result establishes invariance
relative to that model class.

It does not establish invariance across
every conceivable interpretation.

The result depends on the definition of
admissibility and on the expressive and
semantic limitations of the formal language.

A model class may exclude interpretations
that would otherwise challenge the claim.
The adequacy of the model class must therefore
be examined independently.

\subsection{Invariance and the Scope of Formalization}

Not every ontological feature can necessarily
be represented by a single-valued function
or a simple preservation equation.

Some features may be relational, context-dependent,
partially specified, or dependent on structures
that the formal language does not adequately
represent.

In such cases, the basic schema:

\[
    \Phi_R(T(x)) = \Phi_R(x)
\]

may be insufficient.

The formal reconstruction may require:

\begin{enumerate}
    \item relational rather than functional
    representations;

    \item partial functions or restricted domains;

    \item multiple respects of preservation;

    \item context-sensitive interpretation;

    \item higher-order or many-sorted formal
    resources;

    \item a more precise account of the
    transformations being considered.
\end{enumerate}

The introduction of more complex formal
resources must be justified by the requirements
of the subject matter.

Formal complexity is not itself evidence
of philosophical adequacy.

The appropriate formalism is the one that
represents the relevant distinctions with
sufficient precision while avoiding unsupported
ontological commitments.

\subsection{A Protocol for Evaluating Invariance Claims}

Every significant invariance claim in the
formal reconstruction should be examined
according to the following protocol:

\begin{enumerate}
    \item \textbf{Identify the subject:} specify
    the object, structure, state, or relation
    under examination.

    \item \textbf{Specify the domain:} identify
    the domain in which the claim is asserted.

    \item \textbf{Define the transformation:}
    specify the operation or class of operations
    under which preservation is evaluated.

    \item \textbf{Identify the respect:} state
    exactly which feature or relation is
    claimed to be preserved.

    \item \textbf{State the criterion:} define
    what counts as preservation in the relevant
    formal or interpretive setting.

    \item \textbf{Determine the scope:} identify
    whether the claim applies to one
    transformation, a restricted class, or
    a broader class.

    \item \textbf{Examine alternative models:}
    determine whether the result holds across
    the admissible interpretations.

    \item \textbf{Separate formal and ontological
    conclusions:} state which result has been
    established at each level.

    \item \textbf{Examine conceptual relations:}
    identify whether any proposed connection
    to identity, persistence, continuity, or
    stability requires additional premises.

    \item \textbf{Record limitations:} specify
    unresolved assumptions and restrictions.
\end{enumerate}

This protocol makes invariance claims
independently examinable and prevents the
unqualified use of the term from concealing
differences in domain, transformation, or
preservation criterion.

\subsection{Open Problems}

Several problems remain open.

First, the ontology requires a more complete
account of which features can meaningfully
be described as invariant and under what
conditions.

Second, the relation between formal
transformations and ontologically significant
changes requires further justification.

Third, it remains necessary to determine
whether different respects of invariance
can be integrated into a more general
framework without reducing their differences.

Fourth, the relation between invariance
and identity must be examined through
independent arguments rather than assumed
from formal preservation.

Fifth, the possibility of ontological
invariance without a corresponding
single-valued formal invariant remains
an open methodological question.

Sixth, the relation between invariance
and the persistence or continuity of
existents requires a separate analysis
that does not presuppose their equivalence.

These questions identify the limits of
the present formal reconstruction rather
than deficiencies that can be resolved
by notation alone.

\subsection{Concluding Principle}

Invariance is a structured preservation
claim whose meaning depends on the
domain, transformation, respect, criterion,
and scope of its application.

Its formal representation makes the claim
precise enough to be examined across
models and alternative interpretations.
Its ontological significance, however,
depends on the justification of the
interpretation connecting the formal
structure to the intended subject matter.

The governing principle may therefore
be stated as follows:

\begin{quote}
No invariance claim is complete without
specifying what is preserved, under
which transformation, in what respect,
and within which limits. No formal
invariance result, by itself, establishes
the identity, persistence, continuity,
or stability of the subject to which
it is applied.
\end{quote}

This reconstruction establishes the
formal conditions under which invariance
can be discussed without collapsing
distinct ontological questions into
a single notion of preservation.

The next section examines the formal
reconstruction of unity and plurality,
with particular attention to the
distinction between unity, singularity,
aggregation, composition, and plurality
of identities.

% ============================================================
% END OF SECTION 12 OF CHAPTER 07
% ============================================================